\section{Bitcoin-Native Yield}
\label{sec:btc}
% ============================================================

Bitcoin yield historically required custodial wrapping (WBTC~\cite{wbtc}) or centralized lending (BlockFi, Celsius---both insolvent). Babylon~\cite{babylon2024} changed this by enabling native BTC staking through Bitcoin Script timelock covenants, providing economic security to Proof-of-Stake chains without moving BTC off the Bitcoin network.

Teleport composes Bitcoin-native yield through four protocols:

\subsection{Babylon BTC Staking}

Babylon enables BTC holders to stake their BTC to secure PoS chains through a covenant-based locking mechanism on Bitcoin L1~\cite{babylon2024}:

\begin{enumerate}[nosep]
  \item BTC is locked via a self-custodial covenant transaction with a FROST-signed timelock.
  \item The locked BTC provides economic security to participating PoS chains.
  \item Staking rewards accrue in the secured chain's native token.
  \item Unstaking requires a Bitcoin L1 transaction after the timelock expires.
\end{enumerate}

The Teleport bridge vault's BTC is staked through Babylon, with the FROST threshold key serving as both the bridge custody key and the Babylon staking key. This dual use eliminates an additional trust assumption---the same MPC set that secures the bridge also manages the staking position.

\subsection{Liquid BTC Derivatives}

\begin{definition}[BTC Yield Stack]
The Teleport BTC yield stack composes four layers:
\begin{align}
  r_{\text{BTC}} &= r_{\text{Babylon}} + r_{\text{liquid}} + r_{\text{restake}} + r_{\text{DeFi}} \label{eq:btc_yield}\\[4pt]
  &\approx 4.5\% + 1.5\% + 2.0\% + (5\text{--}15\%) \nonumber
\end{align}
\end{definition}

\noindent
Layer composition:

\paragraph{Layer 1: Babylon staking ($r_{\text{Babylon}} \approx 3$--$5\%$).} Raw BTC staking yield from securing PoS chains. The FROST-controlled Taproot address commits to Babylon's staking covenant.

\paragraph{Layer 2: Liquid derivative ($r_{\text{liquid}} \approx 0.5$--$2\%$).} Lombard LBTC~\cite{lombard2024} or SolvBTC~\cite{solvbtc} wraps the Babylon-staked position into a liquid derivative tradeable on EVM chains. The derivative captures the base staking yield plus protocol-specific incentives.

\paragraph{Layer 3: Restaking ($r_{\text{restake}} \approx 1$--$3\%$).} The liquid BTC derivative is restaked through EigenLayer or Symbiotic, providing security to additional AVS protocols in exchange for restaking rewards.

\paragraph{Layer 4: DeFi ($r_{\text{DeFi}} \approx 5$--$15\%$).} yLBTC holders on Lux use their tokens as LP capital, collateral in lending markets, or in yield optimization protocols on the destination chain.

\subsection{CoreDAO Dual Staking}

CoreDAO~\cite{coredao} offers an alternative BTC yield path through dual staking: BTC is delegated to Core validators alongside CORE tokens, earning staking rewards from both networks simultaneously. The yield ($3$--$12\%$ APY) depends on the BTC/CORE delegation ratio and Core network utilization.

\subsection{Security Analysis}

\begin{property}[BTC Custody Invariant]
At all times, the BTC backing yLBTC is:
\begin{enumerate}[nosep]
  \item Held in a FROST 2-of-3 Taproot address (bridge custody), OR
  \item Locked in a Babylon staking covenant with the same FROST key as the spending condition, OR
  \item Represented by a liquid derivative (LBTC, SolvBTC) held by a smart contract controlled by the MPC set.
\end{enumerate}
In all cases, $t$-of-$n$ MPC signers must cooperate to move the BTC. No single party has unilateral withdrawal authority.
\end{property}

% ============================================================
